% Adjoint Method Report
\documentclass[11pt,a4paper]{article}

% --- Packages ---
\usepackage[T1]{fontenc}
\usepackage[utf8]{inputenc}
\usepackage{lmodern}
\usepackage{geometry}
\geometry{margin=1in}
\usepackage{setspace}
\onehalfspacing
\usepackage{microtype}
\usepackage{amsmath,amssymb,amsthm,mathtools}
\usepackage{bm}
\usepackage{bbm}
\usepackage{enumitem}
\usepackage{graphicx}
\usepackage{xcolor}
\usepackage{booktabs}
\usepackage{hyperref}
\usepackage[nameinlink,capitalise]{cleveref}
\usepackage{siunitx}

\numberwithin{equation}{section}
\setlist[itemize]{left=0pt..1em}
\setlist[enumerate]{left=0pt..1em}

% --- Theorem-like ---
\theoremstyle{plain}
\newtheorem{theorem}{Theorem}[section]
\newtheorem{proposition}[theorem]{Proposition}
\newtheorem{lemma}[theorem]{Lemma}
\theoremstyle{definition}
\newtheorem{definition}[theorem]{Definition}
\newtheorem{remark}[theorem]{Remark}

% --- Macros ---
\newcommand{\R}{\mathbb{R}}
\newcommand{\N}{\mathbb{N}}
\newcommand{\ip}[2]{\left\langle #1,\,#2 \right\rangle}
\newcommand{\abs}[1]{\left\lvert #1 \right\rvert}
\newcommand{\norm}[1]{\left\lVert #1 \right\rVert}
\newcommand{\td}{\,\mathrm{d}}
\newcommand{\T}{\mathrm{T}}
\newcommand{\ddt}{\frac{\td}{\td t}}
\DeclareMathOperator*{\argmin}{arg\,min}
\DeclareMathOperator*{\argmax}{arg\,max}

\title{The Adjoint Method: Origins, Derivations, and Practical Use\thanks{Prepared as an extended note on adjoints for optimal control and PDE-constrained optimization.}}
\author{ }
\date{\today}

\begin{document}
\maketitle
\tableofcontents

\begin{abstract}
The adjoint method provides gradients and sensitivities for systems constrained by differential equations with a computational cost that is essentially independent of the dimension of the control or parameter space. This report derives the adjoint equations from first principles in calculus of variations and Lagrangian/KKT theory, presents both ODE- and PDE-based formulations (Pontryagin-type optimality systems), and discusses discrete vs. continuous adjoints, algorithmic differentiation, second-order information via incremental adjoints, and implementation patterns. We include canonical examples (linear-quadratic, elliptic PDE control) and outline how adjoints drive modern gradient-based solvers in large-scale optimization.
\end{abstract}

\section{Motivation and Big Picture}
Consider an optimization problem constrained by a differential equation (ODE or PDE): find controls/parameters $u$ that minimize an objective depending on the corresponding state $y$:
\begin{equation}
\min_{u} \; J(y,u) \quad \text{subject to} \quad C(y,u) = 0.
\label{eq:opt_general}
\end{equation}
Here $C(y,u)=0$ models the governing dynamics (e.g., a PDE) and implicitly defines the parameter-to-state map $u\mapsto y(u)$. The reduced objective is $\hat J(u) := J(y(u),u)$. To use gradient-based optimization efficiently, we need $\nabla \hat J(u)$.

A naive finite-difference gradient costs one forward solve per control component, which is prohibitive when $\dim(u)$ is large. The adjoint method yields $\nabla \hat J(u)$ at a cost comparable to \emph{one} additional solve (the adjoint), independent of $\dim(u)$. This is achieved by introducing a Lagrange multiplier (the adjoint variable) and shifting derivatives via integration by parts or transposition of Jacobians.

\section{Lagrangian View and First-Order Conditions}
Let $Y,U$ be appropriate state and control spaces, and $C:Y\times U\to Z$ encode the constraint. Define the Lagrangian with adjoint $p\in Z$:
\begin{equation}
\mathcal{L}(y,u,p) := J(y,u) + \ip{p}{C(y,u)}.
\end{equation}
Formally, the first-order necessary optimality (KKT) conditions are
\begin{subequations}\label{eq:KKT_general}
\begin{align}
\text{State:}\quad & C(y,u) = 0, \\
\text{Adjoint:}\quad & C_y(y,u)^\ast p + J_y(y,u) = 0, \\
\text{Gradient:}\quad & \nabla \hat J(u) = J_u(y,u) + C_u(y,u)^\ast p = 0 \quad (\text{unconstrained}).
\end{align}
\end{subequations}
If $u$ is constrained (e.g., box or inequality constraints), the gradient condition is replaced by a variational inequality or projection (see \cref{sec:constraints}). The operator $\,\cdot\,^\ast$ denotes the adjoint (transpose in finite dimensions). Equation \eqref{eq:KKT_general} shows how the adjoint $p$ enters the reduced gradient.

\section{Adjoints for ODE-Constrained Problems (Pontryagin)}
\label{sec:ODE}
Consider the optimal control problem
\begin{equation}
\min_{u(\cdot)} \; J(y,u) := \int_0^T \ell(y(t),u(t),t)\,\td t + \phi(y(T))
\end{equation}
subject to the state ODE
\begin{equation}
\dot y(t) = f(y(t),u(t),t), \quad y(0)=y_0.
\end{equation}
Define the Hamiltonian $H(y,u,p,t) := \ell(y,u,t) + p^\T f(y,u,t)$. The Pontryagin minimum principle states that a necessary condition for optimality is the existence of an adjoint $p(t)$ such that
\begin{subequations}\label{eq:PMP}
\begin{align}
\dot y &= f(y,u,t), && y(0)=y_0, \\
-\dot p &= H_y(y,u,p,t) = \ell_y + f_y^\T p, && p(T) = \phi_y\big(y(T)\big), \\
0 &= H_u(y,u,p,t) = \ell_u + f_u^\T p \quad (\text{unconstrained}).
\end{align}
\end{subequations}
For control constraints $u\in \mathcal U$ pointwise, $u(t)$ minimizes $H(\cdot)$ over $\mathcal U$ almost everywhere: $u(t)\in \argmin_{v\in\mathcal U} H(y(t),v,p(t),t)$.

\paragraph{Derivation sketch.} Consider a perturbation $\delta u$ and the corresponding linearized state $\delta y$ solving $\dot{\delta y}=f_y\,\delta y+f_u\,\delta u$, $\delta y(0)=0$. The first variation is
\begin{equation}
\delta J = \int_0^T (\ell_y\,\delta y + \ell_u\,\delta u)\,\td t + \phi_y(y(T))\,\delta y(T).
\end{equation}
Introduce $p$ and integrate by parts:
\begin{align}
\delta J &= \int_0^T (\ell_y\,\delta y + \ell_u\,\delta u)\,\td t + p(T)^\T\delta y(T) - p(0)^\T\delta y(0) - \int_0^T \dot p^\T\delta y\,\td t \\
&\quad - \int_0^T p^\T(f_y\,\delta y + f_u\,\delta u)\,\td t.
\end{align}
Choosing $p$ to satisfy the adjoint ODE in \eqref{eq:PMP} with terminal condition cancels all $\delta y$ terms, leaving
\begin{equation}
\delta J = \int_0^T (\ell_u + f_u^\T p)\,\delta u\,\td t = \int_0^T H_u\,\delta u\,\td t,
\end{equation}
so the $L^2$-gradient of the reduced functional is $g(t)=H_u(y(t),u(t),p(t),t)$.

\section{Adjoints for PDE-Constrained Problems}
\label{sec:PDE}
Let $y$ solve a PDE $A(y,u)=0$ in a spatial domain $\Omega$ with boundary conditions, and consider $J(y,u)$. The Lagrangian is
\begin{equation}
\mathcal{L}(y,u,p) = J(y,u) + \ip{p}{A(y,u)}_{Y^*,Y},
\end{equation}
with $p$ the PDE adjoint. Taking variations and integrating by parts shifts derivatives from $\delta y$ to $p$, producing the \emph{adjoint PDE}. For example, for a linear elliptic control problem
\begin{subequations}\label{eq:elliptic_example}
\begin{align}
-\nabla\cdot(a\nabla y) &= Bu + s && \text{in } \Omega, \\
 y &= 0 && \text{on } \partial\Omega, \\
J(y,u) &= \tfrac{1}{2}\,\norm{y-y_d}_{L^2(\Omega)}^2 + \tfrac{\alpha}{2}\,\norm{u}_{L^2(\Omega)}^2,
\end{align}
\end{subequations}
we obtain the adjoint
\begin{subequations}\label{eq:elliptic_adjoint}
\begin{align}
-\nabla\cdot(a\nabla p) &= y - y_d && \text{in } \Omega, \\
 p &= 0 && \text{on } \partial\Omega,
\end{align}
\end{subequations}
with reduced gradient
\begin{equation}
\nabla \hat J(u) = \alpha u + B^\ast p.
\end{equation}
If $u$ is box-constrained, the optimality condition becomes a projection: $u = \Pi_{[u_{\min},u_{\max}]}\big( -\alpha^{-1} B^\ast p \big)$.

\section{Discrete vs. Continuous Adjoint}
\label{sec:disc_vs_cont}
Two routes exist:
\begin{itemize}
  \item Optimize-then-discretize: derive adjoint equations in the continuum and then discretize them.
  \item Discretize-then-optimize: discretize the state equation and objective, then derive the adjoint of the resulting discrete system.
\end{itemize}
They agree when discretization commutes with differentiation (e.g., Galerkin consistency) and objectives are discretized compatibly. For time-dependent problems, the \emph{discrete adjoint} corresponds to reverse-time propagation through the time-stepping scheme and yields gradients consistent with the discrete objective.

\section{Adjoint and Algorithmic Differentiation}
Reverse-mode algorithmic differentiation (AD) is a discrete adjoint of a computational graph. Backpropagation in deep learning is the adjoint of forward propagation. Many PDE codes can be differentiated via AD frameworks or specialized libraries (e.g., \texttt{dolfin-adjoint}/\texttt{fenics-adjoint} for FEniCS), automatically generating adjoint solves and gradients.

\section{Inequality Constraints and KKT}
\label{sec:constraints}
For control constraints $u\in \mathcal U$ (e.g., box constraints), the stationarity becomes a variational inequality
\begin{equation}
\ip{\nabla \hat J(u), v-u}{} \ge 0 \quad \forall v\in\mathcal U,
\end{equation}
which in Hilbert spaces leads to the projection formula $u=\Pi_{\mathcal U}\big(u-\tau\,\nabla \hat J(u)\big)$ for $\tau>0$. With state inequality constraints (e.g., $g(y)\le 0$), additional multipliers and complementarity conditions arise, and the adjoint gains additional terms $g_y^\ast \mu$ in its right-hand side.

\section{Second-Order Information and Newton-Krylov}
The reduced Hessian action $\nabla^2 \hat J(u)[\delta u]$ can be computed via an \emph{incremental state} and an \emph{incremental adjoint}. Linearize state and adjoint around $(y,u,p)$, solve the incremental systems for $(\delta y, \delta p)$, and assemble the Hessian-vector product. This enables matrix-free Newton-Krylov methods with adjoint-based preconditioners.

\section{Algorithmic Pattern}
A typical gradient-based loop with adjoints:
\begin{enumerate}
  \item Given $u$, solve the state equation for $y$.
  \item Assemble and solve the adjoint equation for $p$ using $(y,u)$.
  \item Form the reduced gradient $g(u)$ (e.g., $g=J_u + C_u^\ast p$).
  \item Update control: projected gradient, L-BFGS, or Newton step with line search.
  \item Stop on stationarity and feasibility criteria.
\end{enumerate}
For transient problems, use checkpointing to manage memory during reverse-time adjoint integration.

\section{Worked Examples}
\subsection{Linear-Quadratic Regulator (discrete time)}
Consider $y_{k+1} = A y_k + B u_k$, $y_0$ given, and cost
\begin{equation}
J = \sum_{k=0}^{N-1} \Big( \tfrac{1}{2} y_k^\T Q y_k + \tfrac{1}{2} u_k^\T R u_k \Big) + \tfrac{1}{2} y_N^\T Q_f y_N.
\end{equation}
The discrete adjoint recursion is
\begin{subequations}
\begin{align}
\lambda_N &= Q_f y_N, \\
\lambda_k &= Q y_k + A^\T \lambda_{k+1}, \quad k=N-1,\dots,0.
\end{align}
\end{subequations}
The reduced gradient is $g_k = R u_k + B^\T \lambda_{k+1}$. Setting $g_k=0$ yields the optimality system; with box constraints, use $u_k=\Pi_{[u_{\min},u_{\max}]}\big(-R^{-1} B^\T \lambda_{k+1}\big)$.

\subsection{Elliptic PDE Control (continuous time)}
The example in \cref{eq:elliptic_example,eq:elliptic_adjoint} leads to a forward elliptic solve for $y$, a backward (same operator) elliptic solve for $p$, and a cheap gradient assembly. This structure is typical: adjoints often reuse the same left-hand side operator as the state (or its transpose), enabling efficient shared preconditioning.

\section{Repository-Specific Mapping (External Mesh Ex3)}
\label{sec:repo_mapping}
We summarize how the adjoint method appears in the provided codebase and how the continuous formulas map to the implemented discrete operators.

\subsection{Model and Spaces}
\begin{itemize}
  \item \textbf{Domain}: $\Omega$ from an external XDMF mesh; control region $\omega \subset \Omega$ is given by a cell tag with integer id $\omega_{\text{id}}$.
  \item \textbf{State/adjoint space} $V$: Lagrange $\mathbb{P}_1$ with homogeneous Dirichlet boundary conditions.
  \item \textbf{Control space} $U$: realized on the same $\mathbb{P}_1$ basis but endowed with the $L^2$ inner product via a mass matrix $M_U$; no boundary conditions are applied in $U$.
\end{itemize}

\subsection{Discrete Operators}
Let $\{\varphi_i\}_{i=1}^n$ be the nodal $\mathbb{P}_1$ basis on $V$. The code assembles
\begin{align}
  A_{ij} &= \int_\Omega \nabla \varphi_i \cdot \nabla \varphi_j \,\td x, &\text{(stiffness)}, \\
  M_{ij} &= \int_\Omega \varphi_i \, \varphi_j \,\td x, &\text{(state mass)}, \\
  (M_U)_{ij} &= \int_\Omega \varphi_i \, \varphi_j \,\td x, &\text{(control mass)}, \\
  B_{ij} &= \int_\omega \varphi_i \, \varphi_j \,\td x. &\text{(coupling on $\omega$)}
\end{align}
Dirichlet boundary conditions are applied to the operator $A$ (rows/columns modified in the assembled FEM operator), and right-hand sides used with $A$ have boundary DOFs zeroed. No boundary treatment is applied to $B$ or $M_U$.

\subsection{State, Adjoint, and Objective}
The discrete state $y \in \R^n$ for a control $u \in \R^n$ solves
\begin{equation}
  A\,y = B\,u \quad (\text{homogeneous Dirichlet via RHS zeroing at boundary DOFs}).
\end{equation}
The target is a pyramid-like profile $y_d(x) = 0.5 - \max\{\lvert x-0.5\rvert, \lvert y-0.5\rvert\}$ interpolated onto $V$. The discrete objective is
\begin{equation}
  F(u) = \tfrac{1}{2}(y-y_d)^\T M (y-y_d) + \tfrac{\beta}{2}\, u^\T M_U \, u.
\end{equation}
The adjoint $p \in \R^n$ solves
\begin{equation}
  A\,p = M\,(y - y_d), \quad \text{with the same Dirichlet treatment as the state.}
\end{equation}

\subsection{Reduced Gradient and Hessian}
With the inner product induced by $M_U$ on the control space, the gradient routine computes
\begin{equation}
  \nabla F(u) = M_U^{-1} B^\T p + \beta\,u.
\end{equation}
This is the Riesz-represented gradient in $U$ such that $\ip{\nabla F(u)}{v}_U = \ip{F'(u)}{v}$ for all $v$. The (matrix-free) Hessian action implemented in the code is
\begin{equation}
  H\,v = M_U^{-1} B^\T A^{-1} M\, A^{-1} B\, v + \beta\, v,
\end{equation}
which is consistent with linearizing the state and adjoint equations around $(y,u,p)$.

\subsection{Algorithmic Details and Verification}
\begin{itemize}
  \item \textbf{Stepsizes:} Gradient descent uses $\alpha=1/L$ with $L$ estimated by a power iteration on $H$; Barzilai--Borwein alternates BB1/BB2 formulas using the $M_U$-induced inner products.
  \item \textbf{Boundary handling:} Apply Dirichlet boundary conditions only to $A$; do not modify $M_U$ or $B$. Zero boundary DOFs on right-hand sides passed to $A$.
  \item \textbf{Consistency check:} Verify gradients via directional derivatives in the $U$ inner product: compare $\frac{F(u+\epsilon v)-F(u)}{\epsilon}$ and $\ip{\nabla F(u)}{v}_U$ for random $v$.
\end{itemize}

\paragraph{Interpretation.} In the continuous setting for $-\Delta y = \chi_\omega u$ with Dirichlet BCs and objective $\tfrac{1}{2}\norm{y-y_d}_{L^2(\Omega)}^2 + \tfrac{\beta}{2}\norm{u}_{L^2(\Omega)}^2$, the adjoint solves $-\Delta p = y - y_d$ with homogeneous Dirichlet boundary conditions, and the reduced gradient is $g = \beta u + B^\ast p$ where $B^\ast$ is the $L^2$ adjoint of $B$ (multiplication by $\chi_\omega$ followed by the appropriate Riesz map). The discrete formulas above implement these relationships faithfully with FEM matrices.

\section{Boundary/Terminal Conditions and Transversality}
For ODEs, terminal costs $\phi(y(T))$ generate terminal conditions $p(T)=\phi_y(y(T))$. For PDEs, integration by parts introduces boundary terms; natural boundary conditions for the adjoint depend on the objective and state BCs. For instance, Dirichlet state BCs typically induce homogeneous Dirichlet adjoint BCs when the objective measures interior misfit; Neumann-type objectives can induce nonzero adjoint Neumann data.

\section{Practical Tips}
\begin{itemize}
  \item \textbf{Consistency:} Match discretization of the objective with the discretization of the forward model to ensure consistent gradients.
  \item \textbf{Scaling:} Nondimensionalize and scale $J$ so state and control terms balance; improves conditioning.
  \item \textbf{Linear solvers:} Reuse factorizations or preconditioners between state and adjoint solves when operators coincide.
  \item \textbf{Checkpointing:} For long transients, use recomputation/checkpointing to trade time for memory during reverse sweeps.
  \item \textbf{Verification:} Check gradients via directional derivatives: $\frac{\hat J(u+\epsilon v)-\hat J(u)}{\epsilon}$ vs. $\ip{\nabla \hat J(u)}{v}{}$.
\end{itemize}

\section{Summary}
The adjoint method is the workhorse for efficient gradient computation in differential equation–constrained optimization. By introducing a Lagrange multiplier and transferring derivatives, adjoints deliver gradients at a cost comparable to one additional solve, enabling scalable optimization, calibration, and control in high-dimensional systems.

\section*{References}
\begin{thebibliography}{9}
\bibitem{Hinze}
M. Hinze, R. Pinnau, M. Ulbrich, S. Ulbrich, \emph{Optimization with PDE Constraints}, Springer, 2009.

\bibitem{Troutman}
J. L. Troutman, \emph{Variational Calculus and Optimal Control}, Springer, 1996.

\bibitem{GilesPierce}
M. B. Giles, N. A. Pierce, ``An introduction to the adjoint approach to design,'' \emph{Flow, Turbulence and Combustion}, 65, 2000.

\bibitem{BorziSchulz}
A. Borz\`{\i}, V. Schulz, \emph{Computational Optimization of Systems Governed by PDEs}, SIAM, 2012.
\end{thebibliography}

\end{document}
